\cotefile{157-1}
\scanpage{257}
\mspage{10}
% archive sheet number: 257
et enfin pour le troisième

\[
\begin{aligned}
&\operatorname{Hom}_{\mathrm{Hot}}(\unsure{\zeta|\xi}\times y,\unsure{\operatorname{Hom}(\mathrm{ex})|\operatorname{Hom}(\mathrm{ext})|\operatorname{Homext}|\operatorname{homext}}(M,N))\\
&\simeq\operatorname{Hom}_{\zeta\times y}(M,N)\simeq\operatorname{Hom}((\zeta\times y)\times M,N),
\end{aligned}
\]

donc toujours la même chose ! \deleted{\unsure{On gagne|On y gagne}} Pas de doutes, ça doit marcher.

S'il faut se donner le $\operatorname{Homext}$ \unsure{de (1)|de 1)} comme donnée supplémentaire, avec de plus un iso. trifonctoriel (12), il faut aussi se donner \unsure{des} iso. trifonctoriels (13), ``\uline{compatibles}'' avec les iso. (12) au sens du calcul précédent ! \deleted{\illegible{1}}

Par exemple, prenant $y={}$type d'homotopie du cercle $S^1$, les formules (13) se spécialisent en

\[
\operatorname{homext}(\Sigma M,N)\simeq
\operatorname{homext}(M,\Omega N)\simeq
\underbrace{\Omega\operatorname{Homext}(M,N)}_{\overset{\mathrm{def}}{=}\operatorname{Hom}_{\mathrm{Hot}}(S^1,-)}.
\tag*{(14)}
\]

Mais cette formule dans $\mathrm{Hot}$ n'est sans doute pas tout à fait la bonne, car le bon formalisme des $\Omega,\Sigma$ se fait dans les catégories \uline{pointées}, et en particulier dans la Catégorie homotopique pointée $\mathrm{Hot}_{\bullet}$.

\hfil$\backslash$\hfil

Je vais \unsure{expliciter} un exemple, en prenant le dérivateur

\[
\mathrm{Der}(A)
\]

\deleted{associé à} ``dérivé'' d'une catégorie abélienne $A$. Donc $\unsure{\mathbb D(e)|D(e)}={}$catégorie dérivée de $A$ --- Pour éviter des
\scanpage{258}
\mspage{11}
% archive sheet number: 258
difficultés \unsure{accessoires|excessives} ici, je vais supposer que $A$ est la cat. des $k$-modules, $k$ un anneau quelconque.

À tout $\unsure{\zeta|\xi}\in\deleted{\unsure{\mathrm{Ob}}}\mathrm{Hot}$ est associé un ``abélianisé''

\[
\zeta_{\mathrm{ab}}\overset{\mathrm{def}}{=}H_{\unsure{\bullet|\cdot}}(\zeta)\in
\underset{\text{ou }D(\mathbb Z\text{-}\mathrm{Mod})}{\underbrace{D(\mathrm{Ab})}}
\quad\begin{aligned}&\text{catégorie dérivée des}\\&\text{groupes abéliens,}\end{aligned}
\]

objet qui n'a de cohomologie qu'en degrés $\leq 0$. Ceci dit, si $M$ est dans $D(k\text{-}\mathrm{Mod})$, on aura

\[
\left\{\begin{aligned}
\zeta\mathbin{\underset{\mathbb D}{\times}}M&\simeq\zeta_{\mathrm{ab}}\overset{\mathbb L}{\otimes}_{\mathbb Z}M\\
\operatorname{hom}_{\mathbb D}(\zeta,M)&\simeq\mathbb R\operatorname{Hom}_{\mathbb Z}(\deleted{\unsure{\zeta}}\zeta_{\mathrm{ab}},M),
\end{aligned}\right.
\tag*{(15)}
\]

\unsure{formules|formule} qui est en \unsure{droit fil} avec la formule de Künneth pour l'homologie

\[
(\zeta\times y)_{\mathrm{ab}}\simeq\zeta_{\mathrm{ab}}\overset{\mathbb L}{\otimes}_{\mathbb Z}y_{\mathrm{ab}},
\tag*{(16)}
\]

ce qui va donner un aspect familier aux formules (9), qui se réduisent :

\[
\left\{\begin{aligned}
\zeta_{\mathrm{ab}}\overset{\mathbb L}{\otimes}_{\mathbb Z}
(y_{\mathrm{ab}}\overset{\mathbb L}{\otimes}_{\mathbb Z}M)
&\simeq(\deleted{\unsure{y}}\zeta_{\mathrm{ab}}\overset{L}{\otimes}_{\mathbb Z}y_{\mathrm{ab}})\overset{L}{\otimes}_{\mathbb Z}M\\
\mathbb R\operatorname{Hom}_{\mathbb Z}(\zeta_{\mathrm{ab}},R\operatorname{Hom}_{\mathbb Z}(y_{\mathrm{ab}},M))
&\simeq R\operatorname{Hom}_{\mathbb Z}(\zeta_{\mathrm{ab}}\overset{L}{\otimes}_{\mathbb Z}y_{\mathrm{ab}},M).
\end{aligned}\right.
\]

% slanted
\margin{expliciter aussi les formules d'adjonction (7) (p.~4).}

D'autre part, existe-t-il bien un $\operatorname{homext}(M,N)$, si $M,N$ \unsure{ds} $D(k\text{-}\mathrm{Mod})$ ? On pense à $\mathbb R\operatorname{Hom}_k(M,N)$, mais pour en faire un type d'homotopie, on doit ``tronquer en degrés $>0$''. On pourrait donc prendre

\[
\operatorname{homext}(M,N)=\tau_{\unsure{\leq|<}0}\bigl(\mathbb R\operatorname{Hom}_k(M,N)\bigr),
\]
\scanpage{259}
\mspage{12}
% archive sheet number: 259
La formule fondamentale (10) est-elle alors valable ? Le deuxième membre est en effet

% slanted
\margin{?? est-ce vrai ???}

\[
\operatorname{Hom}_{\mathrm{Hot}}(\unsure{\zeta|\xi},\deleted{\unsure{\mathbb R\operatorname{Hom}}}H)\simeq
\operatorname{Hom}_{\mathrm{Hotab}}(\zeta_{\mathrm{ab}},H)
\]

car $H$ est un objet abélien de $\mathrm{Hot}$ (cf les topos sur l'abélianisation \ldots), à partir de là c'est gagné :

\[
\begin{aligned}
&\operatorname{Hom}_{\underset{\substack{\approx D(\mathbb Z\text{-}\mathrm{Mod})^{\leq0}\\
\text{sous-cat. }\text{\uline{pleine}}\\\text{de }D(\mathbb Z\text{-}\mathrm{Mod})}}{\underline{\mathrm{Hot}\text{-}\mathrm{ab}}}}
(\zeta_{\mathrm{ab}},\tau_{\leq0}\mathbb R\operatorname{Hom}_k(M,N))
\simeq\operatorname{Hom}_{D(\mathbb Z\text{-}\mathrm{Mod})}
(\zeta_{\mathrm{ab}},\mathbb R\unsure{\operatorname{Hom}}_{\unsure{k}}(M,N))\\
&\hspace{4em}\simeq\operatorname{Hom}_{\deleted{\unsure{D}}(k\text{-}\unsure{M})}(\zeta_{\mathrm{ab}}\times M,N)
\qquad\text{\unsure{oK.}}
\end{aligned}
\]

D'autre part, le \inserted{bi}foncteur

\[
(M,N)\longrightarrow\operatorname{homext}(M,N)=\tau_{\leq0}R\operatorname{Hom}_k(M,N)
\]

se factorise canoniquement par $\mathrm{Hotab}$ et a fortiori par $\mathrm{Hot}$. Le ``bon'' $\operatorname{homext}(M,N)$ serait ici plutôt celui dans $\mathrm{Hotab}$, \unsure{voire|voir} même dans la catégorie $D^{-}(\mathbb Z\text{-}\mathrm{Mod})$, déduite de celle là \deleted{\unsure{par}} ``en rendant le foncteur translation inversible''.

Il est évident d'ailleurs qu'ici le $\zeta_{\mathrm{ab}}\mathbin{\underset{\mathbb D}{\times}}M$, où $\zeta=S^1$, n'est pas la ``bonne'' suspension, qui doit être le décalage \unsure{donc} en guise \deleted{\unsure{de}} $\zeta_{\mathrm{ab}}$ on devrait avoir de la coh. qu'en degré $-1$, et non aussi en degré $0$.
\scanpage{260}
% archive sheet number: 260
\handwritten{diagrammes \unsure{substantiels}}

\handwritten{$\overline{\underline{\mathrm{IV}}}$\unsure{.}}

% printed form
% Mirror image on this face, read from the reverse.
\texttt{GUTS OUTPUT FOR \unsure{PREDONT|PRED0NT} DELIVER AT 0000 - UPDATE AREA}

\texttt{GUTS OUTPUT FOR PREDONT DELIVER AT 0000 - UPDATE AREA}

\texttt{GUTS OUTPUT FOR PREDONT DELIVER AT 0000 - UPDATE AREA}

\texttt{GUTS OUTPUT FOR PREDONT DELIVER AT 0000 - UPDATE AREA}

\texttt{GUTS OUTPUT FOR PREDONT DELIVER AT 0000 - UPDATE AREA}

\texttt{GUTS OUTPUT FOR PREDONT DELIVER AT 0000 - UPDATE AREA}

\texttt{GUTS OUTPUT FOR PREDONT DELIVER AT 0000 - UPDATE AREA}

\texttt{GUTS OUTPUT FOR PREDONT DELIVER AT 0000 - UPDATE AREA}

\texttt{GUTS OUTPUT FOR PREDONT DELIVER AT 0000 - UPDATE AREA}

\texttt{GUTS OUTPUT FOR PREDONT DELIVER AT 0000 - UPDATE AREA}

\texttt{GUTS OUTPUT FOR PREDONT DELIVER AT 0000 - UPDATE AREA}

\texttt{GUTS OUTPUT FOR PREDONT DELIVER AT 0000 - UPDATE AREA}

\texttt{GUTS OUTPUT FOR PREDONT DELIVER AT 0000 - UPDATE AREA}

\texttt{GUTS OUTPUT FOR PREDONT DELIVER AT 0000 - UPDATE AREA}

\texttt{GUTS OUTPUT FOR PREDONT DELIVER AT 0000 - UPDATE AREA}

\texttt{GUTS OUTPUT FOR PREDONT DELIVER AT 0000 - UPDATE AREA}

\texttt{GUTS OUTPUT FOR PREDONT DELIVER AT 0000 - UPDATE AREA}

\texttt{GUTS OUTPUT FOR PREDONT DELIVER AT 0000 - UPDATE AREA}

\texttt{GUTS OUTPUT FOR PREDONT DELIVER AT 0000 - UPDATE AREA}

\texttt{GUTS OUTPUT FOR PREDONT DELIVER AT 0000 - UPDATE AREA}

\texttt{GUTS OUTPUT FOR PREDONT DELIVER AT 0000 - UPDATE AREA}

\texttt{GUTS OUTPUT FOR PREDONT DELIVER AT 0000 - UPDATE AREA}

\texttt{GUTS OUTPUT FOR PREDONT DELIVER AT 0000 - UPDATE AREA}

\texttt{GUTS OUTPUT FOR PREDONT DELIVER AT 0000 - UPDATE AREA}

\texttt{GUTS OUTPUT FOR PREDONT DELIVER AT 0000 - UPDATE AREA}

\texttt{GUTS OUTPUT FOR PREDONT DELIVER AT 0000 - UPDATE AREA}

\texttt{GUTS OUTPUT FOR PREDONT DELIVER AT 0000 - UPDATE AREA}

\texttt{GUTS OUTPUT FOR PREDONT DELIVER AT 0000 - UPDATE AREA}

\texttt{GUTS OUTPUT FOR PREDONT DELIVER AT 0000 - UPDATE AREA}

\texttt{GUTS OUTPUT FOR PREDONT DELIVER AT 0000 - UPDATE AREA}

\texttt{GUTS OUTPUT FOR PREDONT DELIVER AT 0000 - UPDATE AREA}

\texttt{GUTS OUTPUT FOR PREDONT DELIVER AT 0000 - UPDATE AREA}
\scanpage{261}
\mspage{73}
% archive sheet number: 261
\uline{Construction des diagrammes} \uline{substantiels} pour un dérivateur

Soit $\mathbb D$ un prédérivateur, $I$ dans $\mathrm{Dia}$, considérons

\[
\varphi_I:\quad\mathbb D(I)\longrightarrow\underline{\operatorname{Hom}}(I^{\unsure{\mathrm o|\circ}},\mathcal A)
\qquad\text{où }\mathcal A=\unsure{\mathbb D(e)|D(e)}
\]

est la cat. des coefficients fondamentale. Les objets de $\underline{\operatorname{Hom}}(I^{\mathrm o},\mathcal A)$ s'interprètent comme des \uline{diagrammes de type} $I^{\mathrm o}$ dans $\mathcal A$. \inserted{\unsure{Ce langage est surtout commode} si $I$ est \unsure{ordonnée|ordonné}.} Notons, pour $F\in\mathrm{Ob}\,\mathbb D(I)$, par \deleted{\illegible{1}} $(F_i)_{i\in\mathrm{Ob}\,I^{\deleted{\unsure{\mathrm o}}}}$ le $I^{\mathrm o}$-diagramme associé dans $\mathcal A$. Par l'axiome Dér 2, le foncteur $\varphi$

\[
F\longmapsto(F_i)_{i\in\mathrm{Ob}\,I}
\]

est conservatif. Mais il n'a aucune envie d'être fidèle, encore moins pleinement fidèle. Quand on connaît \deleted{\illegible{1}}\unsure{le} diagramme $(F_i)_i$, on est loin de connaître $F$, pas \unsure{$=$|même} à \unsure{isom} non unique près je présume. Il n'est nullement clair que si $F,\deleted{\unsure{E}}\unsure{F'|E'}$ de $\mathbb D(I)$ sont tels que les diagrammes associés soient isomorphes, ils soient eux-mêmes isomorphes.

Pour insister sur ce point, les objets de $\mathbb D(I)$ s'appelleront les ``\uline{diagrammes substantiels} de type $I^{\mathrm o}$ dans $\mathcal A$'' (\unsure{ou} $\mathbb D$-substantiels, s'il y a risque de confusion). Pour étayer l'intuition et simplifier les explications, un objet $F$ de $\mathbb D(I)$ (\inserted{$\mathbb D$-}coefficient sur $I$) sera représenté graphiquement par le diagramme \inserted{``faible'' ou ``\unsure{squelettique}''} correspondant --- sous les regards
\scanpage{262}
\mspage{74}
% archive sheet number: 262
comme une sorte de (\unsure{pâle}) ``squelette'' du diagramme substantiel dont il provient. P. ex. si

\diagram{\begin{tikzcd}[ampersand replacement=\&,column sep=1.6em]
I=\mathbb{\Delta}_0:\quad 0\arrow[r]\&1
\end{tikzcd}}

on note un diagramme substantiel \inserted{\unsure{de|du} type $(\mathbb{\Delta}_{\unsure{1|0}})^{\unsure{\mathrm o|\circ}}$} par le même graphisme, que le \deleted{\illegible{1}} diagramme ordinaire :

\diagram{\begin{tikzcd}[ampersand replacement=\&] F_0\&F_1\arrow[l]\end{tikzcd}\,.}

Si

\diagram{$I={}$\begin{tikzcd}[ampersand replacement=\&,row sep=small,column sep=large]
a\arrow[r]\&c\\
0\arrow[u]\arrow[r]\&b\arrow[u]
\end{tikzcd}}

les diagrammes substantiels de type $I$ se notent

\diagram{\begin{tikzcd}[ampersand replacement=\&,row sep=small,column sep=large]
F_a\arrow[d]\&F_c\arrow[l]\arrow[d]\\
F_0\&F_b\arrow[l]
\end{tikzcd}\,.}

Il faudra bien sûr préciser chaque fois si ce graphisme désigne un diagramme squelettique, ou substantiel. Souvent on part d'un squelette, et on essaye de ``l'étoffer'' avec un peu de chair, pour en faire un diagramme substantiel --- si possible, de façon canonique.

C'est là un point de philosophie dérivatique \unsure{essentielle|essentiel} pour la pratique. On \unsure{procédera|procèdera} souvent de proche en proche, en ``\unsure{substantifiant}'' des sous-\deleted{\unsure{g}}diagrammes correspondant \unsure{à|:} une suite strictement croissante $I_0\subset I_1\ldots\subset I_n=I$ de $I$, au mieux de la situation particulière.
\scanpage{263}
\mspage{75}
% archive sheet number: 263
En voici un exemple typique, en supposant $\mathbb D$ un dérivateur \uline{pointé} (les $\mathbb D(I)$ sont pointées, et les $f^{\unsure{\prime|*}}$ transforment \deleted{\unsure{l'|les}}objets nuls en objets nuls).

Il s'agit de la définition du foncteur \deleted{``\unsure{lacets}''} ``cosuspension'' ou ``\unsure{espaces|espace} de Lacets'' de $\mathcal A$

\[
\Omega:\mathcal A\longrightarrow\mathcal A.
\]

Il est entendu que pour $\deleted{\unsure{N}}M$ de $\mathcal A$, $\Omega\deleted{\unsure{N}}(M)$ doit s'insérer dans un carré commutatif

\diagram{$(C)\qquad$\begin{tikzcd}[ampersand replacement=\&,row sep=large,column sep=large,execute at end picture={
\draw[thin] ([xshift=-1.4em]\tikzcdmatrixname-1-3.west) -- (\tikzcdmatrixname-1-3.north east) -- ([xshift=-1.4em,yshift=.4em]\tikzcdmatrixname-2-3.west) -- (\tikzcdmatrixname-1-3.south east) -- ([xshift=-1.4em]\tikzcdmatrixname-2-3.west) -- (\tikzcdmatrixname-2-3.east);
}]
e\arrow[d]\&\Omega M\arrow[l]\arrow[d]\arrow[r]\&\deleted{e}\arrow[d]\\
M\&e\arrow[l]\arrow[r]\&\deleted{M}
\end{tikzcd}\,,}

où $e$ désigne les éléments unités. Ce carré cependant devrait être substantiel, plus \deleted{\unsure{exactement}}\inserted{précisément}, être un carré substantiel ``\uline{$\mathbb D$-cartésien}''.

Étant entendu qu'un objet de $\mathbb D(I)$ \deleted{\unsure{est|objet}} (carré substantiel de $\mathcal A$) est dit $\mathbb D$-cartésien, s'il est dans l'image essentielle de $\mathbb D(\unsure{\mathbb I_1|I_1})$ par $i_{1*}$, où

\[
i_1:\ \mathbb I_1\lhook\joinrel\longrightarrow I\quad\deleted{\unsure{\mathrm{est}}},
\]

\deleted{\unsure{est défini par l'immersion}} $\mathbb I_1$ étant la \unsure{sous-catég\textsuperscript{ie}|sous-catég\textsuperscript{orie}} ouverte \inserted{de $I$} définie par

\diagram{$\mathbb I_1={}$\begin{tikzcd}[ampersand replacement=\&,row sep=small,column sep=large]
a\&\\
0\arrow[u]\arrow[r]\&b
\end{tikzcd}\,.}
\scanpage{264}
\mspage{76}
% archive sheet number: 264
La \unsure{restr.} du carré \deleted{\unsure{envisagé}} substantiel cherché $(C)$ à la sous-catégorie $\deleted{\unsure{I'}}I_1$ est donc un diagramme substantiel de squelette

\diagram{$(\mathrm{Sq}\qquad$\begin{tikzcd}[ampersand replacement=\&,row sep=large,column sep=large]
e\arrow[d]\&\\
M\&e\arrow[l]
\end{tikzcd}\,,}

dont il se déduit (à isom can. près) en appliquant $i_*$ à ce dernier, de façon à le compléter en

\diagram{\begin{tikzcd}[ampersand replacement=\&,row sep=small,column sep=small]
e\arrow[d]\&\Omega M\arrow[l,dashed]\arrow[d,dashed]\\
M\&e\arrow[l]
\end{tikzcd}\,.}

Il reste donc à décrire une \unsure{substantiation|substantification} canonique du squelette (Sq), en termes de $M$. Cela nous conduit à regarder la sous-catégorie ouverte

\diagram{$I_0=\{0\}: \qquad$\begin{tikzcd}[ampersand replacement=\&,row sep=small,column sep=large]
{}\arrow[r,dotted]\&{}\\
0\,\bullet\arrow[u,dotted]\arrow[r,dotted]\&{}\arrow[u,dotted]
\end{tikzcd}}

de $I_1$ (donc aussi de $I$) réduite à l'objet initial, et l'inclusion

\diagram{\begin{tikzcd}[ampersand replacement=\&] I_0\arrow[r,hook,"i_0"]\&I_1\end{tikzcd}\,,}

et de \deleted{\unsure{déf}} substantifier le squelette Sq par $(i_0)_!(M)$.

On peut dire à \unsure{postériori|post'eriori} que, si on \unsure{connaît}\deleted{\unsure{ssait}} le carré substantiel, il n'y a \inserted{à isom. unique près} qu'une seule façon de l'obtenir comme un $i_{1*}i_{0!}(?)$, avec $?$ dans $\mathbb D(I_0)\simeq\mathcal A$, car $i_{0!}$ et $i_{1*}$ étant pl. fid., il en est de $=$ de leur composé. \deleted{Ainsi le \unsure{fait} qu'il}
\scanpage{265}
\mspage{77}
% archive sheet number: 265
En fait, on voit \deleted{\unsure{par là}} que le foncteur

\[
i_{1*}i_{0!}:\quad\mathcal A\simeq\mathbb D(\deleted{\unsure{?}}I_0)\longrightarrow\mathbb D(I)\deleted{\unsure{\mathrm{sub}}}
\]

induit une équivalence de $\mathcal A$ sur la sous-catégorie pleine de $\mathbb D(I)$ formée des carrés $\mathbb D$-cartésiens \deleted{tels que} dont les fibres en $a$ et $b$ sont nulles. \inserted{(le foncteur quasi-inverse évident est $i^*$, où $i:I_0\lhook\joinrel\longrightarrow I$.)}

Cela signifie donc que dans ce cas, la donnée du squelette \inserted{de nature} très \unsure{particulier|particulière} ($F_a\simeq F_b\simeq0_{\mathcal A}$), plus la contrainte supplémentaire $\mathbb D$-cartésienne, sur le carré cherché, détermine celui-ci à isom unique près en termes de son squelette : il y a une substantiation non seulement ``canonique'' de celui-ci, mais \unsure{$=$.} unique à isom unique près, \deleted{\unsure{en termes}} à cause de la contrainte $\mathbb D$-cartésienne.

Je voudrais commencer par regarder le cas plus simple encore où

\diagram{\begin{tikzcd}[ampersand replacement=\&]
I=\mathbb{\Delta}_1:\quad 0\arrow[r]\&1
\end{tikzcd}\,,}

donc \deleted{\unsure{les diagrammes}} les squelettes envisagés sont simplement \unsure{des|les} flèches

\diagram{\begin{tikzcd}[ampersand replacement=\&] F_1\arrow[r]\&F_0\end{tikzcd}\quad dans $\mathcal A$.}

La philosophie nous oblige, pour les flèches les plus fondamentales dans la \unsure{théorie} du dérivateur $\mathbb D$ (flèches dans $\mathcal A$, ou par extension dans une $\mathbb D(I)\simeq\mathcal A_{\mathbb D_I}$, cat. fondamentale du \deleted{\unsure{$\mathcal A$}} dérivateur \unsure{induit} $\mathbb D_I$), de se demander s'il n'y a pas
\scanpage{266}
\mspage{78}
% archive sheet number: 266
une façon ``canonique'', ``naturelle'', ``voulue par Dieu'', de la ``\unsure{remplumer}'' en un diagramme substantiel.

Les flèches les plus évidentes qui se \deleted{\unsure{présentent}} sont \deleted{\unsure{les}} relatives à un foncteur

\[
f:I\unsure{\text{---}|\longrightarrow}J
\]

de $\mathrm{Dia}$, et aux deux morphismes d'adjonction pour le couple $(f^*,f_*)$, savoir

% slanted
\margin{Bien sûr, problème dual pour $(f_!,f^*)$}

\[
\begin{aligned}
G&\longrightarrow f_*f^*G&&\text{dans }\mathbb D(J)\\
f^*f_*F&\unsure{\text{---}|\longrightarrow}F&&\text{dans }\mathbb D(I).
\end{aligned}
\]

Je vais essayer de remplumer ces deux flèches, et vais généraliser en \deleted{\illegible{2}}\inserted{envisageant la} situation ci-contre, où

\diagram{\begin{tikzcd}[ampersand replacement=\&,row sep=large,column sep=large]
I\arrow[d,"f"']\&F\\
J\&G\arrow[u,"u"']
\end{tikzcd}}

\[
\operatorname{Hom}\unsure{{}^{f}|f|{}_{f}}(G,F)\overset{\mathrm{def}}{=}\operatorname{Hom}(f^*G,F)
\overset{\mathrm{adj}}{\simeq}\operatorname{Hom}(G,\inserted{f_*(F)})
\]

Donc le $=$ $u$ s'interprète par \uline{deux} flèches squelettiques

\[
\left\{\begin{aligned}
f^*(G)&\longrightarrow F&&\text{dans }\mathbb D(I)\\
G&\longrightarrow f_*(F)&&\text{dans }\mathbb D(J).
\end{aligned}\right.
\tag*{(1)}
\]

Les remplumer, c'est les interpréter (\unsure{se|si} faire se peut) en termes de deux objets (si possible canoniques) de $\mathbb D(I\times\mathbb{\Delta}_1)$ et $\mathbb D(J\times\deleted{\unsure{\mathbb D}}\mathbb{\Delta}_1)$ respectivement.

Je vais me guider sur le cas particulier d'un dérivateur trivial

\[
\mathbb D=\mathbb D_{\mathcal M},
\]

où $\mathcal M$ est une catégorie de modèles ``nue'' ($\Sigma=\varnothing$).
\scanpage{267}
\mspage{79}
% archive sheet number: 267
Pour $f$ fixé, on va construire une catégorie auxiliaire $K(f)$, de telle façon que pour \unsure{tte|toute} catégorie $\mathcal M$, les systèmes $(G,F,u)$ d'un $\mathcal M$ préfaisceau \inserted{$G$} sur $J$, d'un autre $F$ sur $I$, et d'\unsure{un|une} $f$-\unsure{hom|flèche} $G\unsure{\text{---}|\longrightarrow}F$, s'identifient (\unsure{fonctoriellement|fonctoriell\textsuperscript{t}}) aux $\mathcal M$-préfaisceaux sur $K(f)$.

Ceci dit, on va essayer d'interpréter les flèches \inserted{(1)} (squelettiques, mais peu importe ici --- tous les \inserted{$\mathbb D$-}squelettes sont \inserted{maintenant} remplumés) en termes de la géométrie de $K(f)$

\[
\left\{\begin{aligned}
&\mathrm{Ob}\,K(f)=\mathrm{Ob}\,K(I)\amalg\mathrm{Ob}\,K(J)\\
&I,\ J\ \text{s'identifient à des sous-catégories pleines de}\\
&\qquad K(f),\ \text{il n'y a pas de flèches d'un }j\in J\text{ en un}\\
&\qquad i\in I.\quad(\text{Donc }I\text{ est ouverte dans }K(f),\ J\ \deleted{\unsure{f}}\ \text{\unsure{son}}\\
&\qquad\text{complémentaire fermé})\\
&\text{Pour }i\in\mathrm{Ob}\,I,\ j\in\mathrm{Ob}\,J,\ \text{on pose}\\
&\qquad\operatorname{Hom}_{\deleted{\unsure{K_J|J}}K(f)}(i,j)=\operatorname{Hom}_J(f(i),j)
\end{aligned}\right.
\]

La composition des flèches dans $K(f)$ est évidente. On va définir une rétraction canonique

\[
\bar f:K(f)\longrightarrow J
\]

de $K(f)$ sur sa sous-catégorie fermée $J$. Donc $\bar f$ est connue sur $\mathrm{Ob}\,J$, et sur $\mathrm{Ob}\,I$ ce sera $f$ :

\[
\left\{\begin{aligned}
\bar f(j)&=j&&\text{si }j\in\mathrm{Ob}\,J\\
\bar f(i)&=f(i)&&\text{si }i\in\mathrm{Ob}\,I
\end{aligned}\right.
\]

\deleted{Pour les flèches} Il reste à définir, pour $x,y\in\mathrm{Ob}\,K(f)$
\scanpage{268}
\mspage{80}
% archive sheet number: 268
\[
\operatorname{Hom}_{\unsure{k(f)|K(f)}}(x,y)\unsure{\text{---}|\longrightarrow}
\operatorname{Hom}_J(\bar f(x),\bar f(y))
\]

Si $x,y\in\mathrm{Ob}\,J$, c'est l'identité, si $x\unsure{,}y\in I$, c'est $\bar f$. Si $x\in I$, $y\in J$, on doit définir une flèche

\[
\operatorname{Hom}(f(x),y)\longrightarrow\operatorname{Hom}(f(x),y)
\]

et on prend encore l'identité. \deleted{Le foncteur} On a

\diagram{\begin{tikzcd}[ampersand replacement=\&,column sep=small,row sep=large]
J\arrow[r,hook,"j"]\arrow[dr,equal]\&K(f)\arrow[d,"\bar f"]\& I\arrow[l,hook,no head,"i"']\\
\&J\&
\end{tikzcd}}

\[
\left\{\begin{aligned}
\bar f\circ j&=\mathrm{id}_J\\
j\circ\bar f&\overset{\deleted{\unsure{?}}}{=}
\begin{cases}\deleted{\unsure{*}}j&\text{sur }J\\jf&\text{sur }I\end{cases}
\end{aligned}\right.
\qquad\text{i.e.}\quad
\underset{\deleted{\text{sur }I}}{\boxed{\begin{aligned}
j\bar f j&=j\ \text{sur }J\\
j\bar f i&=jf\ \text{sur }I
\end{aligned}}}
\]

et on a \deleted{\illegible{1}}

\[
\begin{aligned}
\alpha:\mathrm{id}_K&\longrightarrow j\bar f&&\text{i.e. hom fonctoriel}\\
x&\overset{\alpha(\unsure{x|0})}{\longmapsto}j(\bar f(x))&&x\in\mathrm{Ob}\,K
\end{aligned}
\]

défini ainsi :

\[
\left\{\begin{aligned}
&\text{si }x\in J,\quad\alpha(x)=\mathrm{id}_x,\ \text{\unsure{s.à|s}}:\quad
 x\xrightarrow{\sim}j(\bar f(x))=j(x)=\deleted{\unsure{j}}x\\
&\text{si }x\in I\qquad\alpha(x)\mathrel{\unsure{\doteq|=}} x\underset{\text{ds }K(x)}{\longrightarrow}f(x),\ \text{ou }
 \underbrace{f(x)\longrightarrow f(x)}_{\text{ds }I}\\
&\hspace{18em}\alpha(x)=\mathrm{id}_{f(x)},
\end{aligned}\right.
\]

On doit vérifier la \deleted{\unsure{compatibilité}}\inserted{fonctorialité} pour $x$ variable dans $K$\ldots

Soit $\Phi$ un $\mathcal M$-préfaisceau sur $K$. Il définit

\[
\begin{aligned}
F&\mathrel{\unsure{\doteq|=}}i^*(\Phi)&&\in\mathrm{Ob}\,\mathbb D_{\mathcal M}(I)\\
G&=j^*(\Phi)&&\in\mathrm{Ob}\,\mathbb D_{\mathcal M}(J)
\end{aligned}
\]

D'autre part $f^*(G)$ s'interprète comme $(j\bar f)^*(\Phi)|I$ car

\[
\begin{aligned}
f^*(G)(i)&=G(\underset{\bar f(i)\,\unsure{=|\parallel}\ \in J}{\underline{f(i)}})
 =j^*(\Phi)(\bar f(i))=\deleted{\unsure{G(f(i))}}\Phi(j\bar f(i))\\
&=(j\bar f)^*(\Phi)(i)
\end{aligned}
\]
\scanpage{269}
\mspage{81}
% archive sheet number: 269
Donc

\[
f^*(G)\simeq(j\bar f)^*(\Phi)|I
\]

Or la flèche

\[
\alpha:\mathrm{id}_K\longrightarrow j\bar f,\quad
\deleted{\unsure{i\longrightarrow j\bar f(i)}}
\]

donne

\[
\underset{\substack{\parallel\\\Phi}}{\mathrm{id}_K^*(\Phi)}
\xleftarrow{\Phi(\alpha)}(j\bar f)^*(\Phi)
\qquad\text{fonctoriel en }\Phi
\]

d'où, en prenant la restriction à $I$

\[
i^*(\Phi)\overset{\mathrm{def}}{=}F\xleftarrow{u}f^*G.
\]

En fait, sans tellement se tortiller, on devrait pouvoir montrer que \deleted{\unsure{en}} pour \unsure{toute} catégorie $\mathcal M$, la catégorie $\underline{\operatorname{Hom}}(K^{\unsure{\mathrm o|\circ}},\mathcal M)$ des $\mathcal M$-préfaisceaux sur $K$ est \deleted{équivalente}\inserted{isomorphe} à celle des triples $(F,G,u)$, avec $F\in\mathrm{Ob}\,\underline{\operatorname{Hom}}(I^{\mathrm o},\mathcal M)$, $G\in\underline{\operatorname{Hom}}(J^{\mathrm o},\mathcal M)$, et $u:f^*(G)\longrightarrow F$. \deleted{\unsure{Mais}} \unsure{Ainsi}, $K$ apparaît comme solution d'un pb \deleted{\illegible{1}} universel

\deleted{\unsure{rel. au foncteur covariant}}

\[
\underline{\operatorname{Hom}}(K^{\mathrm o},\mathcal M)\longrightarrow
\{F,G,u\ \text{comme dessus}\}(\mathcal M)
\]

Mais pour y préciser la flèche $\longrightarrow$ de façon \unsure{conceptuelle|conceptuellement} simple, il est bon de se donner $K$ avec les structures supplémentaires

\diagram{\begin{tikzcd}[ampersand replacement=\&,row sep=tiny,column sep=small]
J\arrow[r,hook,"j"]\&K\arrow[l,bend left=35,"\bar f"]\&I\arrow[l,hook,no head,"i"']\\
\&\bar f j=\mathrm{id}_J\&
\end{tikzcd}
\quad$,\quad\mathrm{id}_K\xrightarrow{\alpha}j\bar f$.}

Et ce système
\scanpage{270}
\mspage{82}
% archive sheet number: 270
$\{K,i,j,\bar f,\alpha\}$ est (sauf erreur) défini à iso. can. près par la condition de représenter le foncteur $\mathcal M\longmapsto\{F,G,u\}^{\unsure{\mathrm o|\circ}}(\mathcal M^{\mathrm o})$\ldots

En particulier, le système défini par $G\in\mathrm{Ob}\,\mathbb D_{\mathcal M}(J)$

\[
(f^*(\deleted{\unsure{F}}G),\deleted{\unsure{f|F}}G,\mathrm{id}:f^*(G)\longrightarrow f^*(G))
\]

est défini sauf erreur par $\bar f^*(G)$ sur $K(f)$. On se demande de trouver une interprétation similaire \deleted{\unsure{pour}}

\[
(F,f_*(F),f^*f_*F\longrightarrow F)
\]

($F\in\mathrm{Ob}\,\mathbb D(I)$). Le candidat naturel serait ici $\Phi=i_*(F)$. Admettons ce point (i.e. \inserted{surtout} la vérification du fait que $j^*(i_*(F))\simeq f_*(F)$).

Tout ceci montre que si $\mathbb D$ est à nouveau un dérivateur quelconque, on peut, par les formules précédentes, trouver un objet canonique \inserted{$\Phi$} de $\mathbb D(K)$, en termes soit d'un \inserted{$\mathbb D$-coeff} $F$ sur $I$ (en prenant $\Phi=i_{\unsure{1|*}}(F)$), soit d'un $\mathbb D$-coeff: $G$ sur $J$ (en prenant $\Phi=\bar f^*(\unsure{G})$.

D'autre part, on peut espérer définir canoniquement

\[
\begin{aligned}
\mathbb D(K)\deleted{\unsure{(f)}}&\longrightarrow\mathbb D(I\times\mathbb{\Delta}_1)\\
\mathbb D(K)&\longrightarrow\mathbb D(J\times\mathbb{\Delta}_1)
\end{aligned}
\]

dont les composés avec les foncteurs ``squelette''

\[
\begin{aligned}
\mathbb D(I\times\mathbb{\Delta}_1)&\longrightarrow\underline{\operatorname{Hom}}(\mathbb{\Delta}_1^{\mathrm o},\mathbb D(I))\\
\mathbb D(J\times\mathbb{\Delta}_1)&\unsure{\text{---}|\longrightarrow}\underline{\operatorname{Hom}}(\mathbb{\Delta}_1^{\mathrm o},\mathbb D(J))
\end{aligned}
\]
\scanpage{271}
\mspage{83}
% archive sheet number: 271
\unsure{redonnent} \inserted{(pour $\Phi\in\mathrm{Ob}\,\mathbb D(\unsure{K|k})$)} les flèches squelettiques

\[
\left\{\begin{aligned}
f^*(G)&\xrightarrow{u}F\\
G&\xrightarrow{v}f_*F
\end{aligned}\right.
\]

associées à $(F,G,u)$. De façon précise, $f^*(G)\longrightarrow F$ \deleted{s'identifie} se définit en termes de $\mathbb D$, et $\Phi\in\mathrm{Ob}\,\mathbb D(K)$, comme ci-dessus, et $v$ s'en déduit par adjonction.

L'idée naturelle ici, serait de ramener le pb de l'étoffement des flèches d'adjonction squelettiques, au cas des \unsure{imm.|immersions} ouvertes ou fermées. Notons que les sous-cat. ouvertes de $K$ correspondent justement aux foncteurs $K\longrightarrow\mathbb{\Delta}_{\unsure{9|1|q}}$, en prenant les images inverses de l'ouvert générique. Il y a donc un unique

\[
\unsure{\varphi|\phi}:K(f)\unsure{\text{---}|\longrightarrow}\mathbb{\Delta}_1
\]

avec $I=\varphi^{-1}(\{0\})$. Donc
